\documentclass[12pt]{article}
\usepackage[spanish]{babel}
\usepackage[T1]{fontenc}
\usepackage[utf8]{inputenc}
\usepackage{amsmath,amssymb}
\usepackage{amsthm}
\usepackage{geometry}
\geometry{margin=2.5cm}
\usepackage{comment}
\renewcommand{\familydefault}{\sfdefault} % texto en sans serif
% \usepackage{lmodern} % alternativa si prefieres serif moderna
\usepackage{mathrsfs}

% Matemáticas y entornos
\usepackage{amsmath,amssymb}
% (amsthm no es obligatorio, pero útil si quieres \begin{proof})
\usepackage{amsthm}
\newtheorem{definicion}{Definición}[section]
\newtheorem{lema}{Lema}[section]
\newtheorem{proposicion}{Proposición}[section]
\newtheorem{teorema}{Teorema}[section]
\newtheorem{corolario}{Corolario}[section]
\newtheorem{observacion}{Observación}[section]
\newtheorem{ejemplo}{Ejemplo}[section]
\newtheorem{conjetura}{Conjetura}[section]
\newtheorem{problema}{Problema abierto}[section]

% Figuras y tablas
\usepackage{graphicx}
\usepackage{float}       % para [H]
\usepackage{subcaption}
\usepackage{booktabs,tabularx,array}

% Listas
\usepackage{enumitem}

% Títulos (opcional; quítalo si no lo usas o si cambias a wiley-article)
\usepackage{titlesec}

% Bibliografía con natbib (autor-año)
\usepackage[authoryear]{natbib}
\bibliographystyle{plainnat}

% Hipervínculos (SIEMPRE al final)
\usepackage[hidelinks]{hyperref}

\title{Control Informacional\\
\large Proyecto de Investigación II --- Introducción}

\author{
\textbf{Alejandro Cruz López} \\
Universidad Autónoma Metropolitana \\
Unidad Iztapalapa
}

\date{
\vspace{1em}
\textbf{Asesores} \\[0.3em]
Dr.\ Asael Fabián Martínez Martínez \\
Dr.\ Joaquín Delgado Fernández \\
\vspace{1em}
25 de mayo de 2026
}

\begin{document}
\maketitle

\input{CI_ProyectoII_Intro_base.tex}

\section*{Introducción}

Este documento fija el punto de partida del Proyecto de Investigación II.
La meta es estudiar transformaciones sobre el símplex de probabilidades finitas
cuando dos mecanismos actúan de manera conjunta: la evolución tipo Markov y la
actualización tipo Bayes. El interés principal no es solo describir cada
operador por separado, sino entender qué ocurre cuando se aplican en distinto
orden, cuánto cambia el resultado y cómo medir esa diferencia de forma
geométrica e informacional.

La estrategia de trabajo será gradual. Primero se consolidará el marco ya
validado en \texttt{ProInv\_I.tex}, que proporciona la base estructural sobre el
símplex, las correcciones multiplicativas y la geometría informacional. Después
se tratará \texttt{CI.tex} como un repositorio exploratorio de ideas, sin
asumir todavía que sus proposiciones sean definitivas. Finalmente,
\texttt{CI\_ProyectoII.tex} servirá como el documento activo del proyecto,
pero cada resultado deberá quedar sustentado con referencias claras y con
demostraciones revisadas antes de considerarse válido.

En esta etapa, el objetivo central es formular con precisión el problema de
conmutatividad entre Markov y Bayes y, si es posible, obtener cotas que
cuantifiquen el defecto de ordenamiento. Una cota de este tipo no solo diría si
dos operaciones conmutan, sino qué tan costoso es intercambiarlas. Eso la
convierte en una herramienta útil para análisis de robustez, filtrado,
asimilación de datos y control secuencial.

\begin{observacion}
El papel de la bibliografía será decisivo. Ninguna afirmación se tomará como
definitiva mientras no esté respaldada por referencias verificables o por una
demostración interna completa. En este sentido, el trabajo combinará lectura
crítica, depuración matemática y construcción progresiva del esquema final.
\end{observacion}

\section*{Objetivo general}

Construir un marco matemático en el que la interacción entre actualización
bayesiana y evolución markoviana pueda describirse, compararse y cuantificarse
de forma precisa sobre el símplex de probabilidad.

\section*{Líneas de avance}

\begin{itemize}[nosep]
  \item Consolidar la base teórica validada.
  \item Revisar cuidadosamente las referencias bibliográficas.
  \item Reescribir o depurar las proposiciones que aún sean exploratorias.
  \item Formular y analizar la cota de no conmutatividad Bayes-Markov.
  \item Traducir esa cota a interpretaciones concretas en control, filtrado y
        actualización secuencial.
\end{itemize}

\section*{Enunciados propuestos}

Los siguientes enunciados están formulados como un bloque de trabajo para
desarrollar el proyecto con rigor. La intención no es adelantar demostraciones,
pero sí fijar con precisión el lenguaje formal mínimo que hará falta para la
parte central de la investigación.

\begin{definicion}[Símplex de probabilidad e interior]
Sea $n\geq 2$. El \emph{símplex de probabilidad} es
\[
\Delta_n:=\Bigl\{X\in\mathbb{R}^n:\ X_i\geq 0\ \text{para todo } i,\ 
\sum_{i=1}^n X_i=1\Bigr\}.
\]
Su interior relativo es
\[
\Delta_n^\circ:=\{X\in\Delta_n:\ X_i>0\ \text{para todo } i\}.
\]
\end{definicion}

\begin{definicion}[Matriz estocástica por filas]
Una matriz $K\in\mathbb{R}^{n\times n}$ se llama \emph{estocástica por filas}
si satisface
\[
K_{ij}\geq 0
\quad\text{y}\quad
\sum_{j=1}^n K_{ij}=1
\qquad\text{para todo } i\in\{1,\dots,n\}.
\]
\end{definicion}

\begin{definicion}[Acciones admisibles]
Fijados parámetros $R>0$ y $0<\varepsilon\leq 1/n$, se definen:
\[
\mathcal{K}_\varepsilon
:=
\Bigl\{K\in\mathbb{R}^{n\times n}:\ K \text{ es estocástica por filas y } K_{ij}\geq \varepsilon\ \forall i,j\Bigr\},
\]
\[
\mathcal{U}_R
:=
\bigl\{v\in\mathbb{R}^n:\ \|v\|_\infty\leq R\bigr\},
\]
y
\[
\mathcal{E}_R
:=
\bigl\{L\in\mathbb{R}_{>0}^n:\ \log L\in\mathcal{U}_R\bigr\}
=
\bigl\{L\in\mathbb{R}_{>0}^n:\ e^{-R}\leq L_i\leq e^{R}\ \forall i\bigr\}.
\]
\end{definicion}

\begin{definicion}[Operador bayesiano]
Para $L\in\mathbb{R}_{>0}^n$ y $X\in\Delta_n^\circ$, el \emph{operador bayesiano}
asociado a $L$ es
\[
B_L(X):=\frac{X\circ L}{\langle X,L\rangle},
\qquad
\langle X,L\rangle:=\sum_{i=1}^n X_iL_i.
\]
En coordenadas,
\[
\bigl(B_L(X)\bigr)_i=\frac{X_iL_i}{\sum_{j=1}^n X_jL_j},
\qquad i=1,\dots,n.
\]
\end{definicion}

\begin{definicion}[Defecto de ordenamiento]
Sean $K\in\mathcal{K}_\varepsilon$ y $L\in\mathcal{E}_R$. El \emph{defecto de
ordenamiento} en el estado $X\in\Delta_n^\circ$ es
\[
\mathcal{D}_{K,L}(X):=B_L(XK)-B_L(X)K.
\]
\end{definicion}

\begin{definicion}[Política admisible de horizonte $T$]
Fijados un horizonte $T\geq 1$, un estado inicial $X_0\in\Delta_n^\circ$, un
objetivo $P\in\Delta_n^\circ$ y un parámetro $\lambda>0$, una \emph{política
admisible} es una sucesión $\pi=\{(a_t,A_t)\}_{t=0}^{T-1}$ tal que:
\begin{itemize}[nosep]
  \item $a_t\in\{M,B\}$ para todo $t$;
  \item si $a_t=M$, entonces $A_t\in\mathcal{K}_\varepsilon$;
  \item si $a_t=B$, entonces $A_t\in\mathcal{E}_R$;
  \item el estado evoluciona por
  \[
  X_{t+1}=
  \begin{cases}
  X_tA_t & \text{si } a_t=M,\\[0.2em]
  B_{A_t}(X_t) & \text{si } a_t=B.
  \end{cases}
  \]
\end{itemize}
El costo asociado a $\pi$ es
\[
J_T(X_0;\pi)
:=
D_{\mathrm{KL}}(P\|X_T)
\lambda\sum_{t=0}^{T-1} c_{a_t}(A_t,X_t),
\]
donde
\[
c_M(K,X):=D_{\mathrm{KL}}(XK\|X),
\qquad
c_B(L,X):=D_{\mathrm{KL}}(B_L(X)\|X).
\]
La función de valor óptimo se define por
\[
V_T(X_0):=\inf_{\pi}J_T(X_0;\pi),
\]
donde el ínfimo se toma sobre todas las políticas admisibles de horizonte $T$.
\end{definicion}

\begin{lema}[Compacidad de las acciones admisibles]
El conjunto $\mathcal{K}_\varepsilon$ es un poliedro compacto y convexo. El
conjunto $\mathcal{E}_R$ es un cubo compacto y convexo. En particular, ambos
conjuntos son no vacíos si y solo si $0<\varepsilon\leq 1/n$ y $R>0$.
\end{lema}

\begin{lema}[Cotas elementales de los operadores]
Sea $K\in\mathcal{K}_\varepsilon$ y $X\in\Delta_n$. Entonces
\[
(XK)_j\geq \varepsilon
\qquad\text{para todo } j,
\]
y en particular $XK\in\Delta_n^\circ$. Sea ahora $L\in\mathcal{E}_R$ y
$X\in\Delta_n^\circ$. Entonces
\[
e^{-2R}X_i\leq \bigl(B_L(X)\bigr)_i\leq e^{2R}X_i
\qquad\text{para todo } i,
\]
y, en particular, $B_L(X)\in\Delta_n^\circ$.
\end{lema}

\begin{corolario}[Trayectorias admisibles en el interior]
Toda trayectoria generada a partir de un estado inicial de $\Delta_n^\circ$ por
una política admisible de horizonte finito permanece en $\Delta_n^\circ$ en
todos sus tiempos intermedios.
\end{corolario}

\begin{proposicion}[Regularidad del defecto]
La aplicación
\[
\mathcal{D}_{K,L}:\Delta_n^\circ\longrightarrow \mathbb{R}^n
\]
es de clase $C^\infty$ en $X$. Además,
\[
\sum_{i=1}^n \mathcal{D}_{K,L}(X)_i=0,
\]
de modo que $\mathcal{D}_{K,L}(X)\in T_X\Delta_n^\circ$ para todo
$X\in\Delta_n^\circ$.
\end{proposicion}

\begin{teorema}[Existencia de políticas óptimas de horizonte finito]
Para todo $T\geq 1$, todo $X_0\in\Delta_n^\circ$, todo $P\in\Delta_n^\circ$,
todo $\lambda>0$, todo $R>0$ y todo $0<\varepsilon\leq 1/n$, el problema de
optimización definido por $V_T(X_0)$ admite al menos una política óptima.
\end{teorema}

\begin{corolario}[Continuidad de la función de valor]
Para cada horizonte fijo $T$, la función de valor óptimo $V_T$ es continua en
$\Delta_n^\circ$.
\end{corolario}

\begin{teorema}[Cota local de suboptimalidad por defecto]
Sea $\mathcal{C}\subset\Delta_n^\circ$ un compacto. Entonces existe una
constante $C_{\mathcal{C},P}>0$ tal que, siempre que
\[
B_L(XK),\ B_L(X)K\in\mathcal{C},
\]
se tiene
\[
\bigl|D_{\mathrm{KL}}(P\|B_L(XK)) - D_{\mathrm{KL}}(P\|B_L(X)K)\bigr|
\leq
C_{\mathcal{C},P}\,\|\mathcal{D}_{K,L}(X)\|_1.
\]
\end{teorema}

\begin{corolario}[Invarianza por defecto nulo]
Si $\mathcal{D}_{K,L}(X)=0$, entonces
\[
D_{\mathrm{KL}}(P\|B_L(XK))=D_{\mathrm{KL}}(P\|B_L(X)K).
\]
Más generalmente, cualquier funcional terminal continuo que dependa solo del
estado final es invariante bajo el intercambio de orden cuando el defecto se
anula.
\end{corolario}

\section*{Resultados preliminares por ruta}

Los enunciados siguientes se ordenan de menor a mayor ambición. Los primeros
son puentes técnicos desde la base validada; los últimos son metas abiertas que
todavía requieren demostración o refinamiento. La idea no es cerrar el
problema, sino fijar con rigor qué piezas hacen falta para llegar a un punto no
trivial.

\subsection*{Ruta 0: Puente desde la base validada}

\begin{proposicion}[Bayes como automorfismo del interior]\label{prop:bayes-automorfismo}
Sea $L\in\mathbb R_{>0}^n$. Entonces la aplicación
$B_L:\Delta_n^\circ\to\Delta_n^\circ$ es de clase $C^\infty$. Si se define
$L^{-1}:=(L_1^{-1},\dots,L_n^{-1})$, entonces
\[
B_{L^{-1}}\circ B_L=\operatorname{id}_{\Delta_n^\circ}
\qquad\text{y}\qquad
B_L\circ B_{L^{-1}}=\operatorname{id}_{\Delta_n^\circ}.
\]
\end{proposicion}

\begin{proposicion}[Ley de composición bayesiana]\label{prop:bayes-composicion}
Para cualesquiera $L,M\in\mathbb R_{>0}^n$ se cumple
\[
B_L\circ B_M=B_{L\circ M},
\qquad
L\circ M:=(L_1M_1,\dots,L_nM_n).
\]
En particular, $B_{\alpha L}=B_L$ para todo $\alpha>0$, y la familia de
operadores bayesianos es conmutativa bajo composición.
\end{proposicion}

\begin{teorema}[Gradiente Fisher de la divergencia KL]\label{teo:gradiente-kl}
Para cada $P\in\Delta_n^\circ$, el funcional
$E_P(X):=D_{\mathrm{KL}}(P\|X)$ tiene gradiente riemanniano respecto de la
métrica de Fisher dado por
\[
\operatorname{grad}_g E_P(X)=X-P.
\]
En consecuencia, el flujo de gradiente satisface
\[
\dot X=P-X
\]
y su solución explícita es
\[
X(t)=P+e^{-t}(X_0-P).
\]
\end{teorema}

\begin{corolario}[Descenso estricto y mínimo global]\label{cor:minimo-kl}
A lo largo del flujo Fisher--KL, el valor de $E_P$ no aumenta y es
estrictamente decreciente fuera del equilibrio. Además, $P$ es el único
minimizador global de $E_P$ en $\Delta_n^\circ$.
\end{corolario}

\begin{proposicion}[Lipschitz local de las transiciones]\label{prop:lipschitz-transiciones}
Sea $C\subset\Delta_n^\circ$ un compacto. Entonces existen constantes
$L_M(C,\varepsilon)>0$ y $L_B(C,R)>0$ tales que, para todo $X,Y\in C$, todo
$K\in\mathcal K_\varepsilon$ y todo $L\in\mathcal E_R$,
\[
\|XK-YK\|_1\le L_M(C,\varepsilon)\|X-Y\|_1,
\qquad
\|B_L(X)-B_L(Y)\|_1\le L_B(C,R)\|X-Y\|_1.
\]
\end{proposicion}

\subsection*{Ruta 1: Alcanzabilidad restringida}

\begin{definicion}[Conjunto alcanzable restringido]
Dado $X_0\in\Delta_n^\circ$ y un horizonte $t\ge 0$, el \emph{conjunto
alcanzable restringido} $\mathcal R_t(X_0)$ es el conjunto de todos los
estados finales que pueden obtenerse a partir de $X_0$ mediante una sucesión de
$t$ acciones, donde cada paso es o bien una acción de Markov $X\mapsto XK$ con
$K\in\mathcal K_\varepsilon$, o bien una acción bayesiana $X\mapsto B_L(X)$ con
$L\in\mathcal E_R$.
\end{definicion}

\begin{definicion}[Conjunto alcanzable cuantizado]
Fijado $M\ge 1$, el \emph{conjunto alcanzable cuantizado} $\mathcal R_t^{(M)}(X_0)$
es el subconjunto de $\mathcal R_t(X_0)$ formado por las trayectorias cuyos
estados intermedios pertenecen al retículo
\[
\Delta_n^{(M)}:=\{X\in\Delta_n:\ MX_i\in\mathbb Z_{\ge 0}\ \text{para todo }i\}.
\]
\end{definicion}

\begin{proposicion}[Recursión del conjunto alcanzable]\label{prop:recursion-alcanzable}
Para todo $t\ge 0$ y todo $X_0\in\Delta_n^\circ$,
\[
\mathcal R_{t+1}(X_0)
=
\bigl\{YK:\ Y\in\mathcal R_t(X_0),\ K\in\mathcal K_\varepsilon\bigr\}
\cup
\bigl\{B_L(Y):\ Y\in\mathcal R_t(X_0),\ L\in\mathcal E_R\bigr\}.
\]
\end{proposicion}

\begin{teorema}[Compacidad del conjunto alcanzable]\label{teo:compacticidad-alcanzable}
Para todo horizonte finito $t\ge 0$ y todo $X_0\in\Delta_n^\circ$, el conjunto
$\mathcal R_t(X_0)$ es no vacío y compacto en $\Delta_n^\circ$.
\end{teorema}

\begin{corolario}[Persistencia uniforme del interior]\label{cor:interior-persistente}
Para cada $t$ finito y cada $X_0\in\Delta_n^\circ$, existe $\delta_t>0$ tal
que
\[
\mathcal R_t(X_0)\subset \Delta_{n,\delta_t},
\qquad
\Delta_{n,\delta}:=\{X\in\Delta_n:\ X_i\ge\delta\ \text{para todo }i\}.
\]
\end{corolario}

\begin{teorema}[Inalcanzabilidad robusta en tiempo finito]\label{teo:inalcanzabilidad-robusta}
Para todo $t\ge 0$ y todo $X_0\in\Delta_n^\circ$, existe $P\in\Delta_n^\circ$
tal que
\[
P\notin\mathcal R_t(X_0)
\quad\text{y}\quad
\operatorname{dist}_1\!\bigl(P,\mathcal R_t(X_0)\bigr)>0.
\]
En particular, la alcanzabilidad restringida es no trivial en todo horizonte
finito.
\end{teorema}

\begin{conjetura}[Pérdida por cuantización de orden $O(1/M)$]\label{conj:cuantizacion}
Para $M$ suficientemente grande, la pérdida de alcanzabilidad debida a la
cuantización satisface una cota del tipo
\[
\inf_{Q\in\mathcal R_t^{(M)}(X_0)}D_{\mathrm{KL}}(P\|Q)\le \frac{C}{M},
\]
para una constante $C>0$ que depende del horizonte y de las restricciones, pero
no de $M$.
\end{conjetura}

\begin{problema}[Caracterización geométrica de $\mathcal R_t(X_0)$]\label{prob:caracterizacion-reach}
Determinar condiciones necesarias y suficientes sobre $X_0$, $P$,
$\mathcal K_\varepsilon$, $\mathcal E_R$ y $t$ para que $P\in\mathcal R_t(X_0)$.
En particular, describir la geometría de $\mathcal R_t(X_0)$ dentro de
$\Delta_n^\circ$.
\end{problema}

\subsection*{Ruta 2: Control óptimo informacional}

\begin{proposicion}[Compacidad del espacio de políticas]\label{prop:politicas-compactas}
Fijado un horizonte $T$, el espacio de políticas admisibles de horizonte $T$
es compacto como unión finita de productos compactos. Además, el funcional
$\pi\mapsto J_T(X_0;\pi)$ es continuo sobre dicho espacio.
\end{proposicion}

\begin{teorema}[Existencia de política óptima de horizonte finito]\label{teo:existencia-politica-optima}
Para todo $T\ge 1$, todo $X_0\in\Delta_n^\circ$, todo $P\in\Delta_n^\circ$,
todo $\lambda>0$, todo $R>0$ y todo $0<\varepsilon\le 1/n$, el problema de
optimización definido por $V_T(X_0)$ admite al menos una política óptima.
\end{teorema}

\begin{proposicion}[Ecuación de Bellman restringida]\label{prop:bellman}
La función de valor satisface la recursión hacia atrás
\[
\begin{aligned}
V_\tau(X)
&=
D_{\mathrm{KL}}(P\|X)
+\min\!\Bigl(
\inf_{K\in\mathcal K_\varepsilon}
\bigl[\lambda\,c_M(K,X)+V_{\tau-1}(XK)\bigr],\\
&\hspace{4.2em}
\inf_{\log L\in\mathcal U_R}
\\
&\hspace{6.2em}
\bigl[\lambda\,c_B(L,X)+V_{\tau-1}(B_L(X))\bigr]
\Bigr).
\end{aligned}
\]
Además, bajo las restricciones $(\mathcal K_\varepsilon,\mathcal U_R)$, el
ínfimo no se anula de manera trivial para $\tau$ finito.
\end{proposicion}

\begin{corolario}[Cota de referencia para el valor óptimo]\label{cor:cota-referencia}
Para todo horizonte $T$ y todo $X\in\Delta_n^\circ$ se cumple
\[
0\le V_T(X)\le D_{\mathrm{KL}}(P\|X).
\]
\end{corolario}

\begin{proposicion}[Monotonicidad con respecto al horizonte]\label{prop:monotonia-horizonte}
Para todo $X\in\Delta_n^\circ$ y todo $T\ge 0$,
\[
V_{T+1}(X)\le V_T(X).
\]
En particular, la sucesión $\{V_T(X)\}_{T\ge 0}$ es decreciente y está
acotada inferiormente por $0$.
\end{proposicion}

\begin{corolario}[Límite asintótico de la función de valor]\label{cor:valor-limite}
Para cada $X\in\Delta_n^\circ$, existe el límite
\[
V_\infty(X):=\lim_{T\to\infty}V_T(X),
\]
y satisface $0\le V_\infty(X)\le D_{\mathrm{KL}}(P\|X)$.
\end{corolario}

\begin{corolario}[Separación del objetivo]\label{cor:valor-cero}
Para todo $T\ge 0$, se tiene $V_T(P)=0$. Además, si $X\neq P$, entonces
$V_T(X)>0$.
\end{corolario}

\begin{conjetura}[Política de tipo \emph{bang-bang} para $n=2$]\label{conj:bangbang}
Para $n=2$ y restricciones simétricas (mismos parámetros $\varepsilon$ y $R$
para ambas acciones), existe una partición de $\Delta_2^\circ$ en dos regiones
$\mathcal M$ y $\mathcal B$ tal que la política óptima es $a^*=M$ en
$\mathcal M$ y $a^*=B$ en $\mathcal B$. En particular, la política óptima es de
acción pura en cada región y no requiere alternancia.
\end{conjetura}

\begin{conjetura}[Ventaja de la alternancia para $n\ge 3$]\label{conj:alternancia}
Para $n\ge 3$, existen configuraciones $(X_0,P,\mathcal K_\varepsilon,
\mathcal U_R)$ para las cuales
\[
J_T^*(X_0)<\min\!\bigl(J_T^M(X_0),\,J_T^B(X_0)\bigr),
\]
donde $J_T^*$ es el costo de la política óptima y $J_T^M$, $J_T^B$ son los
costos de las políticas puras (solo Markov y solo Bayes, respectivamente). Es
decir, la alternancia de operadores produce un costo estrictamente menor que
cualquier política pura.
\end{conjetura}

\begin{conjetura}[Límite continuo igual al flujo Fisher--KL]\label{conj:limite-continuo}
Para $T$ fijo y $\Delta t=T/N$ con costo de acción $c_M(K,X)$, en el límite
$N\to\infty$ la trayectoria de la política óptima converge a la solución de la
ecuación diferencial
\[
\dot X=P-X,
\]
que es el flujo de gradiente Fisher--KL de $D_{\mathrm{KL}}(P\|\cdot)$.
\end{conjetura}

\subsection*{Ruta 3: Defecto de ordenamiento}

\begin{proposicion}[Regularidad conjunta del defecto]\label{prop:defecto-suave}
La aplicación
\[
\mathcal D:\Delta_n^\circ\times\mathcal K_\varepsilon\times\mathcal E_R
\longrightarrow \mathbb R^n,
\qquad
(X,K,L)\longmapsto \mathcal D_{K,L}(X),
\]
es de clase $C^\infty$. En particular, para cada compacto $C\subset\Delta_n^\circ$
existe una constante $M_C>0$ tal que
\[
\|\mathcal D_{K,L}(X)-\mathcal D_{K,L}(Y)\|_1\le M_C\|X-Y\|_1
\]
para todo $X,Y\in C$, todo $K\in\mathcal K_\varepsilon$ y todo $L\in\mathcal E_R$.
\end{proposicion}

\begin{proposicion}[Criterios triviales de anulación]\label{prop:defecto-trivial}
Si $K=I_n$ o si $L=c\,\mathbf 1$ para alguna constante $c>0$, entonces
\[
\mathcal D_{K,L}(X)=0
\qquad\text{para todo }X\in\Delta_n^\circ.
\]
\end{proposicion}

\begin{proposicion}[Lipschitz local de la divergencia terminal]\label{prop:lipschitz-terminal}
Sea $C\subset\Delta_n^\circ$ un compacto y sea $P\in\Delta_n^\circ$ fijo.
Entonces existe una constante $G_{C,P}>0$ tal que, para todo $U,V\in C$,
\[
\bigl|D_{\mathrm{KL}}(P\|U)-D_{\mathrm{KL}}(P\|V)\bigr|
\le G_{C,P}\,\|U-V\|_1.
\]
\end{proposicion}

\begin{problema}[Caracterización algebraica de $\mathcal D_{K,L}\equiv 0$]\label{prob:conmutacion}
Determinar condiciones necesarias y suficientes sobre $K$ y $L$ para que
\[
\mathcal D_{K,L}(X)=0
\quad\text{para todo }X\in\Delta_n^\circ.
\]
\end{problema}

\begin{proposicion}[Subvariedad de conmutación]\label{prop:subvariedad}
Sea
\[
\mathcal M_{K,L}:=\bigl\{X\in\Delta_n^\circ:\ \mathcal D_{K,L}(X)=0\bigr\}.
\]
Si en un punto $X_0\in\mathcal M_{K,L}$ el diferencial
$D_X\mathcal D_{K,L}(X_0):T_{X_0}\Delta_n^\circ\to T_{X_0}\Delta_n^\circ$
tiene rango constante $r$, entonces $\mathcal M_{K,L}$ es, localmente cerca de
$X_0$, una subvariedad inmersa de codimensión $r$.
\end{proposicion}

\begin{teorema}[Cota local de suboptimalidad por defecto]\label{teo:cota-defecto}
Sea $\mathcal C\subset\Delta_n^\circ$ un compacto. Entonces existe una
constante $C_{\mathcal C,P}>0$ tal que, siempre que
\[
B_L(XK),\ B_L(X)K\in\mathcal C,
\]
se tiene
\[
\bigl|D_{\mathrm{KL}}(P\|B_L(XK)) - D_{\mathrm{KL}}(P\|B_L(X)K)\bigr|
\le
C_{\mathcal C,P}\,\|\mathcal D_{K,L}(X)\|_1.
\]
\end{teorema}

\begin{corolario}[Invariancia por defecto nulo]
Si $\mathcal{D}_{K,L}(X)=0$, entonces
\[
D_{\mathrm{KL}}(P\|B_L(XK))=D_{\mathrm{KL}}(P\|B_L(X)K).
\]
Más generalmente, cualquier funcional terminal continuo que dependa solo del
estado final es invariante bajo el intercambio de orden cuando el defecto se
anula.
\end{corolario}

\begin{conjetura}[Cota global uniforme de suboptimalidad]\label{conj:cota-global}
Existe una constante $C=C(\varepsilon,R)>0$ tal que, para toda
$K\in\mathcal K_\varepsilon$, toda $L$ con $\log L\in\mathcal U_R$ y todo
$X\in\Delta_n^\circ$,
\[
\bigl|J(X;K,L)-J(X;L,K)\bigr|
\le
C\cdot\|\mathcal D_{K,L}(X)\|_1,
\]
donde $J(X;K,L):=D_{\mathrm{KL}}(P\|B_L(XK))$ y
$J(X;L,K):=D_{\mathrm{KL}}(P\|B_L(X)K)$.
\end{conjetura}

\end{document}
